\chapter{Introduzione}
\label{ch:introduzione}

% BOZZA — testo di partenza dagli appunti, da rivedere con la relatrice.

Il software testing è una componente critica del processo di sviluppo: una
suite di test di alta qualità permette di individuare incongruenze tra le
specifiche di un sistema e la sua implementazione, e di catturare regressioni
introdotte durante l'evoluzione del codice. Scrivere buoni test è tuttavia
un'attività costosa e ripetitiva, spesso trascurata in tutto o in parte.

I modelli linguistici di grandi dimensioni (LLM) hanno mostrato capacità
notevoli nella generazione di codice, ma i risultati nella generazione di
test per progetti reali restano deludenti: nel benchmark TestGenEval il
miglior modello raggiunge una copertura media del codice di appena il 35\%,
principalmente per la difficoltà dei modelli nel ragionare sull'esecuzione
del codice.

\section{Obiettivo della tesi}
% TODO: rifinire con la relatrice.
Questa tesi indaga se e quanto la copertura dei test generati da un LLM
migliori fornendo al modello: (i) contesto strutturale sul codice sotto test
(ad esempio l'albero sintattico astratto), (ii) strumenti software di
supporto, e (iii) un ciclo di feedback iterativo basato sui risultati
dell'esecuzione dei test e sul report di coverage. L'idea centrale è
sostituire la capacità --- carente --- del modello di \emph{prevedere} il
comportamento del codice con l'\emph{osservazione} dei risultati reali
dell'esecuzione.

\section{Struttura del documento}
% TODO: completare quando i capitoli saranno stabili.
Il Capitolo~\ref{ch:stato_arte} inquadra lo stato dell'arte sulla
generazione automatica di test. Il Capitolo~\ref{ch:core} descrive il
sistema proposto. Il Capitolo~\ref{ch:risultati} presenta gli esperimenti e
i risultati. Seguono conclusioni e sviluppi futuri.
